\writeSectionFrame{\presentationTitle}{Das Problem}
\frameWithImageOnly{\BILDERPFAD/problem1.jpg}
\note{
	Die abstrakte Fabrik bietet erheblich mehr Möglichkeiten als die einfache
	Fabrikmethode, ist aber natürlich auch komplexer. Es geht um das Erstellen
	von ganzen Produktfamilien.
}

\begin{frame}{Unser Möbelhaus}
	\begin{center}
 		\includegraphics[scale=0.45]{\BILDERPFAD/problem2.png}
 	\end{center}
\end{frame}

\note{
 	Wir entwerfen einen Simulator für ein Möbelhaus. Dort stellen wir eine Produktfamilie
 	im viktorianischen Stil vor. Eine Produktfamilie besteht aus verschiedenen Produkten wie
 	z.B. Sofa, Sessel, Stühle, Tisch usw. Die Produkte einer Familie sind in demselben Stil,
 	z.B. im viktorianischen oder modernen Stil. 
}

\begin{frame}{Unser Möbelhaus}
	\begin{center}
 		\includegraphics[scale=0.35]{\BILDERPFAD/problem3.png}
 	\end{center}
\end{frame}

\note{
 	Ein Kunde hat nun Sessel und Sofa im viktorianischen Stil gekauft und anschließend noch einen
 	Stuhl dazu bestellt. Wir haben allerdings einen absolut unpassenden Stuhl geliefert. Das sollte
 	man natürlich vermeiden. Der Stuhl müsste auf jeden Fall aus derselben Produktfamilie kommen.
 	Wir müssen also Produktfamilien bauen, bei denen alle Produkte möglichst ähnlich aussehen und
 	zueinander passen. 
}

\begin{frame}{Unser Möbelhaus}
	\begin{center}
 		\includegraphics[scale=0.35]{\BILDERPFAD/problem4.png}
 	\end{center}
\end{frame}

\note{
 	Zudem müssen neue Produktfamilien leicht ins System integrierbar sein, ohne dass wir dafür viel
 	existierenden Code ändern müssen. 
}